\section{Background}
\label{sec:background}

This section reviews the source language, the functional representation produced by Aeneas, and the target logic used by \texttt{Rust2Isabelle}.

\subsection{Rust}

Rust is a systems programming language designed to provide memory safety
without garbage collection. Its type system assigns each value a unique
owner and controls access through borrowing.
A shared reference \texttt{\&T} permits read-only access and may be duplicated, whereas a mutable reference \texttt{\&mut T} permits exclusive access.
The borrow checker uses lifetimes to ensure that references do not outlive their owners and that mutable and shared access do not overlap in unsafe ways.
For safe Rust, these rules prevent errors such as use-after-free and unsynchronized data races without requiring a runtime memory manager~\cite{jung2017rustbelt}.

Rust programs also make extensive use of algebraic data types, pattern matching, traits, parametric polymorphism, and const generics.
These features are relevant to functional verification because they affect both the types of translated functions and the evidence passed to generic code.
Mutable borrows present the main semantic difficulty: an update through \texttt{\&mut T} must eventually be reflected in the value owned by the caller.
Rust's compiler represents checked programs using the MIR, a typed control-flow graph~\cite{rustc-mir} intended for compiler analyses and optimization rather than interactive proofs.

\subsection{Aeneas's Pure AST}

Aeneas takes as input LLBC, the typed and structured representation of Rust MIR produced by Charon~\cite{ho2025charon}. It symbolically executes LLBC using an ownership-centric semantics.
For its supported safe-Rust fragment, this process removes places, loans, and lifetime annotations and produces a side-effect-free functional representation called the Pure AST~\cite{ho2022aeneas}.
Shared borrows become ordinary values. A mutable borrow may produce a forward value together with a backward function.
After the borrowed value is updated, the backward function reconstructs the corresponding owner.
This representation captures the effect of mutation without carrying an explicit heap into the theorem prover.

The Pure AST also makes failures explicit. Operations that may panic or fail, such as checked arithmetic and bounds-checked indexing, return a result value.
Its declarations include types, functions, globals, trait declarations and implementations, const-generic parameters, and recursive groups.
Aeneas orders these declarations by dependency before passing them to a backend.
\texttt{Rust2Isabelle} starts at this boundary. It does not repeat Rust borrow checking or the symbolic execution from LLBC to Pure.

\subsection{Isabelle/HOL}

Isabelle is a generic interactive theorem prover that supports different
object logics.
Isabelle/HOL instantiates it with classical higher-order logic and is the logic used in this work~\cite{nipkow2002isabelle}.
An Isabelle theory contains named declarations and proofs and may import other theories.
Isabelle/HOL provides polymorphic types, datatypes, records, higher-order functions, and support for recursive and inductive definitions.
Its proof environment also offers automation tools such as Sledgehammer~\cite{bohme2010sledgehammer}.

Several properties of Isabelle/HOL affect the translation of the Pure AST.
Functions are total, and every HOL type is nonempty.
Declarations use commands such as \texttt{datatype}, \texttt{record}, \texttt{definition}, and \texttt{function}, and referenced types and constants must already be in scope.
These constraints shape the representation and translation choices presented in the next section.
